\documentclass[10pt,a4paper]{amsart}
\usepackage[margin=19mm]{geometry}
\usepackage{amsmath,amssymb,mathtools}
\usepackage[scr=rsfs,cal=euler]{mathalfa}
\usepackage{xfrac}
\usepackage[unicode,hidelinks,bookmarks=false]{hyperref}
\newcommand{\ie}{i.e.\ }
\newcommand{\cf}{cf.\ }
\newcommand{\resp}{respectively\ }
\newcommand{\Subsection}{\subsection}
\providecommand{\nobreakdash}{\nobreak\hbox{-}}
\setlength{\emergencystretch}{2em}
\multlinegap0pt
\usepackage{url}
\usepackage{booktabs}
\usepackage{nicefrac}
\usepackage{bm}
\usepackage{multirow}
\usepackage{float}
\newtheorem{theorem}{Theorem}
\newtheorem{assumption}{Assumption}
\newtheorem{remark}{Remark}
\newcommand{\Mtwo}{\mathcal{M}_2}
\newcommand{\DM}{\mathcal{D}_M}
\newcommand{\lcoh}{\lambda_{\mathrm{coh}}}
\newcommand{\lvis}{\lambda_{\mathrm{vis}}}
\newcommand{\Lop}{\mathcal{L}}
\newcommand{\Pioff}{\Pi_{\mathrm{off}}}
\newcommand{\LC}{\mathrm{LC}}
\newcommand{\Tr}{\operatorname{Tr}}
\newcommand{\Dcohlay}{\Delta_{\mathrm{coh}}^{\mathrm{lay}}}
\newcommand{\weff}{\widehat{\omega}}
\newcommand{\aeff}{\widehat{a}}
\newcommand{\eqdef}{\mathrel{\vcenter{\baselineskip0.5ex \lineskiplimit0pt
  \hbox{\scriptsize.}\hbox{\scriptsize.}}}=}
\providecommand{\ket}[1]{\lvert #1\rangle}
\providecommand{\bra}[1]{\langle #1\rvert}
\begin{document}
\title[A Coherence Law for Trainability in Noisy Equivariant Quantu]{A Coherence Law for Trainability in Noisy Equivariant Quantum Neural Networks}
\author{Hassan Ugail; Centre for Visual Computing and Intelligent Systems; United Kingdom; Newton Howard; United States; Hassan Ugail, Newton Howard}
\date{}
\maketitle
\noindent\emph{Source and licence.} A Coherence Law for Trainability in Noisy Equivariant Quantum Neural Networks. Version arXiv:2606.30688v1 (CC BY 4.0). This material is distributed under the Creative Commons Attribution 4.0 International licence, http://creativecommons.org/licenses/by/4.0/, as recorded in the arXiv OAI licence metadata for this exact version and shown on the abstract page https://arxiv.org/abs/2606.30688v1. Full rights notice: licenses/arxiv-2606.30688-CC-BY-4.0.txt.\par\smallskip
\noindent\textit{Editorial scope.} Counted unit: the original \texttt{theorem} environment \texttt{thm:perturb} at source lines 206--224 together with its complete proof at lines 242--246; the delivered document keeps source lines 82--247 so that every referenced label and every citation resolves inside it. Native editable source is the paper's own concatenated TeX; the AMS wrapper is editorial formatting only. No publisher class, upstream script or unrelated artwork is redistributed.
\section{Methods}\label{sec:methods}

\subsection{Model, readout and gradient observables}\label{sec:model}

We study a $U(1)$-equivariant quantum neural network on a cycle graph $C_n$ with periodic boundary conditions. The conserved charge is $\hat N=\sum_j(I-Z_j)/2$, and the circuit acts on a fixed input that lies in a single charge sector $r$, with $P_r$ the orthogonal projector onto that sector. The input is the localised computational-basis state $\rho_0=\ket{\psi_0}\bra{\psi_0}$ with $\ket{\psi_0}=\ket{1^r 0^{n-r}}$. The variational circuit is a depth-$L$ brickwork ansatz built from two gate families that commute with $\hat N$ \cite{nguyen2024,schatzki2024,larocca2022groupinv}. Each layer applies trainable single-qubit rotations $R_z(\theta_{\ell,j})$ on every site, followed by nearest-neighbour $XY$ hopping gates $\exp[-i\beta_{\ell,e}(X_{e_1}X_{e_2}+Y_{e_1}Y_{e_2})]$ on alternating edges of the cycle. The hopping parameters are fixed and treated as part of the architecture, so the only trainable degrees of freedom are the single-qubit rotation angles.

This study examines how noise degrades the parameter gradients of a fixed-state equivariant response rather than learning from a distribution of encoded data. The state the circuit processes is held fixed, and the object of interest is the second moment of the readout gradient with respect to the trainable angles. This is the quantity whose survival determines whether the circuit can be optimised at all, and it is well defined independently of any data-encoding map. The intra-sector off-diagonal coherence that carries the gradient is generated dynamically by the $XY$ hopping gates, which mix the in-sector basis states as the excitation propagates through the variational layers. The readout-visible coherence studied throughout the paper is therefore a property of the variational evolution acting on the fixed input, and the sector-supported state remains in the same charge sector throughout the noiseless circuit.

The output is the sector-projected local readout,
\begin{equation}
    O_B = P_r Z_0 P_r ,
    \label{eq:readout}
\end{equation}
and the noisy response is,
\begin{equation}
    f^\gamma_\theta=\Tr\!\bigl[O_B\,\Phi^\gamma_\theta(\rho_0)\bigr],
    \label{eq:prediction}
\end{equation}
where $\Phi^\gamma_\theta$ denotes the layered circuit channel with single-qubit Markovian noise of rate $\gamma$ applied once per layer \cite{lindblad1976,gks1976,breuerpetruccione2002}. Trainability is quantified by the prediction-gradient second moment,
\begin{equation}
    \Mtwo(\theta_i;\gamma)
    =
    \mathbb{E}_{\theta}
    \Bigl[\bigl(\partial_{\theta_i}f^\gamma_\theta\bigr)^2\Bigr],
    \label{eq:m2pred}
\end{equation}
with expectations taken over a small-box parameter prior of width $\delta_{\mathrm{init}}$ \cite{grant2019} and derivatives evaluated by the parameter-shift rule \cite{mitarai2018,schuld2019,wierichs2022}. The central response variable is the relative degradation,
\begin{equation}
    \DM(\gamma)
    =
    \log\!\bigl[\Mtwo(\theta_i;0)\,/\,\Mtwo(\theta_i;\gamma)\bigr],
    \label{eq:dm}
\end{equation}
which measures on a logarithmic scale how much active gradient strength the noise has removed.

\subsection{Active gradients and the backward light cone}\label{sec:lightcone}

Before asking how noise damages a gradient, it helps to ask where a gradient can live at all. The parameter-shift identity expresses a derivative as a commutator response,
\begin{equation}
    \partial_{\theta_i}f^0_\theta
    =
    -\tfrac{i}{2}\,\Tr\!\bigl[O_B(\ell_i)\,[Z_{q_i},\rho^\theta_{\ell_i}]\bigr],
    \label{eq:psr}
\end{equation}
where $O_B(\ell_i)$ is the readout evolved backwards in the Heisenberg picture. Each brickwork layer propagates operator support by at most one bond, so $O_B(\ell_i)$ is supported on the backward light cone $\LC(O_B)$ of the readout qubit, and outside that cone the commutator vanishes by causal locality. Because every gate commutes with $\hat N$, the trace restricts further to the charge sectors within the cone. A theorem makes this precise, subject to two conditions on the fixed background.

\begin{assumption}[Light-cone algebra non-degeneracy]\label{ass:nondegen}
For the fixed hopping background $\beta^\dagger$, the commutator between the light-cone-evolved readout operator and the active generator is not identically zero on the sector-restricted light-cone subspace.
\end{assumption}

\begin{assumption}[Light-cone sector weight]\label{ass:weight}
The fixed input places a non-vanishing charge weight inside the readout light cone, uniformly in system size. There is a constant $w_0>0$, independent of $n$ at fixed $L$ and $r$, such that the light-cone reduced state $\rho_{\LC}$ satisfies $\Tr[P^{\LC}_{q\geq1}\,\rho_{\LC}]\geq w_0$, where for the single-excitation sector $P^{\LC}_{q\geq1}=P^{\LC}_{q=1}$ projects onto the one-excitation subspace of the cone. The condition excludes inputs in which the excitation delocalises so strongly that its weight inside a fixed-size cone vanishes as $1/n$.
\end{assumption}

\begin{theorem}[Light-cone reduction of active gradients]\label{thm:locality}
Let $\LC(O_B)$ be the backward light cone of the readout, let $|\LC|$ be the number of qubits in it, and let $d_{\LC}=\sum_{q=0}^{r}\binom{|\LC|}{q}$ be the sector-restricted light-cone dimension. Active gradients are supported only inside the cone, since $\partial_{\theta_i}f^0_\theta=0$ whenever $\theta_i\notin\LC(O_B)$. Under Assumptions~\ref{ass:nondegen} and~\ref{ass:weight}, for any active parameter $\theta_i\in\LC(O_B)$ and the small-box prior of width $\delta_{\mathrm{init}}$,
\begin{equation}
    \Mtwo(\theta_i;0)
    \geq
    c_{\LC}\bigl(L,r,\delta_{\mathrm{init}},\beta^\dagger,O_B,w_0\bigr)
    >0 ,
    \label{eq:lcbound}
\end{equation}
where $c_{\LC}$ is independent of the total system size $n$ at fixed $L$ and $r$.
\end{theorem}

The locality statement that gradients vanish outside the cone needs no assumption on the input. The $n$-independent lower bound is the stronger claim and rests on Assumption~\ref{ass:weight}, which guarantees that the excitation keeps finite weight inside the fixed-size cone as $n$ grows. The theorem frames everything that follows. Only the light-cone parameters carry signal, so the question becomes what noise does to the gradient amplitude inside the cone, and that amplitude is carried by intra-sector coherence.

\begin{proof}[Proof of Theorem~\ref{thm:locality}]
We adopt the generator convention $R_z(\theta)=\exp(-i\theta Z/2)$, so the half-angle generator is $Z/2$ and the parameter-shift rule uses shifts of $\pi/2$. If parameter $\theta_i$ acts on qubit $q_i$ in layer $\ell_i$, the parameter-shift identity gives Equation~\eqref{eq:psr}. Parameters for which the commutator $[Z_{q_i},O_B(\ell_i)]$ vanishes identically on the sector-restricted Hilbert space are classified as inactive and excluded from Theorem~\ref{thm:locality}.

The proof begins with the sector-restricted light-cone reduction. Each brickwork layer propagates operator support by at most one bond, so the Heisenberg-evolved readout $O_B(\ell_i)$ is supported on the backward light cone $\LC(O_B)$. Outside the cone the commutator vanishes by causal locality, and because all gates commute with $\hat N$ the trace further restricts to the charge sectors $q\leq r$ within the cone, whose dimension is $d_{\LC}=\sum_{q=0}^{r}\binom{|\LC|}{q}$, replacing the full sector dimension $\binom{n}{r}$.

Next, the prior factorises over inactive coordinates. Writing $\theta=(\theta^{\LC},\theta^{\bar\LC})$ for parameters inside and outside the cone, the derivative for an active parameter depends only on $\theta^{\LC}$. The small-box prior is a product measure, so the coordinates outside the cone integrate to unit mass and
\begin{equation}
    \mathbb{E}_\theta\!\bigl[|\partial_{\theta_i}f^0_\theta|^2\bigr]
    =
    \int_{[-\delta_{\mathrm{init}},\delta_{\mathrm{init}}]^{p_{\LC}}}
    |\partial_{\theta_i}f^0_{\theta^{\LC}}|^2\,
    \frac{\mathrm{d}^{p_{\LC}}\theta^{\LC}}{(2\delta_{\mathrm{init}})^{p_{\LC}}},
    \label{eq:factorisation}
\end{equation}
where $p_{\LC}$, the number of light-cone parameters, depends on the local circuit geometry but not on the total number of qubits at fixed $L$ and $r$. By Assumption~\ref{ass:nondegen} the integrand is a continuous function of $\theta^{\LC}$ that is not identically zero, so there is an open set $U_{\LC}$ on which $|\partial_{\theta_i}f^0_{\theta^{\LC}}|^2\geq\mu(\beta^\dagger)>0$. The in-cone signal $\mu(\beta^\dagger)$ is proportional to the active-sector light-cone weight $\Tr[P^{\LC}_{q\geq1}\rho_{\LC}]$, so Assumption~\ref{ass:weight} keeps it from vanishing as $1/n$ when the excitation delocalises. Combining the factorisation with the non-degeneracy and sector-weight floor gives
\begin{equation}
    \Mtwo(\theta_i;0)
    \geq
    \frac{\mathrm{vol}(U_{\LC})}{(2\delta_{\mathrm{init}})^{p_{\LC}}}\,
    w_0\,\widetilde{\mu}(\beta^\dagger)
    \eqdef
    c_{\LC}(L,r,\delta_{\mathrm{init}},\beta^\dagger,O_B,w_0)>0,
    \label{eq:lowerbound}
\end{equation}
a constant independent of $n$ at fixed $L$ and $r$, where $\widetilde{\mu}$ is the in-cone signal normalised by the sector weight. The bound is conditional on the fixed non-degenerate background $\beta^\dagger$ and on Assumption~\ref{ass:weight}, and is not claimed to be uniform over degenerate backgrounds or maximally delocalised inputs.
\end{proof}

\subsection{The readout-visible aligned coherence rate}\label{sec:rate}

To make the coherence dependence quantitative, we need an operator-level rate. Let $\Phi^\gamma=I+\gamma\Lop+O(\gamma^2)$ be the small-noise expansion of the single-qubit channel and let $\Pioff$ project onto operators on the off-diagonal block of sector $r$ in the eigenbasis of $\hat N$. The worst-case coherence-contraction rate is,
\begin{equation}
    \lcoh
    \eqdef
    \sup_{\substack{A=\Pioff A\\ \|A\|_F=1}}
    -\mathrm{Re}\,\langle A,\Pioff\Lop(A)\rangle_F ,
    \label{eq:lcohdef}
\end{equation}
the largest decay rate of the restricted generator on the intra-sector off-diagonal subspace in the Frobenius inner product. This worst-case rate is convenient but not fundamental, because it maximises over all off-diagonal directions while the gradient occupies only one. The active gradient of parameter $\theta_i$ is generated by the commutator direction $G_i(\theta)=\Pioff[Z_{q_i},O_B(\ell_i;\theta)]$, the readout-visible mode through which the projected readout actually responds. The fundamental object is the readout-visible aligned rate, the ensemble-weighted decay rate of the restricted generator along this direction,
\begin{equation}
    \lvis(\theta_i)
    \eqdef
    \mathbb{E}_{\theta}\!\left[
    \frac{-\mathrm{Re}\,\langle G_i(\theta),\Pioff\Lop(G_i(\theta))\rangle_F}
         {\langle G_i(\theta),G_i(\theta)\rangle_F}
    \right].
    \label{eq:lvisdef}
\end{equation}
By construction, the aligned rate is bounded above by the worst-case rate, since for every $\theta$ the supremum in Equation~\eqref{eq:lcohdef} runs over a set that contains the normalised gradient direction, so
\begin{equation}
    \lvis(\theta_i)\leq\lcoh
    \label{eq:visbound}
\end{equation}
for every active parameter. The two coincide when the restricted generator acts isotropically on the off-diagonal block, because every normalised direction then decays at the common rate. The distinction is invisible for an isotropic family but decisive for structured noise, where a channel can contract directions orthogonal to $G_i$ and so carry a large $\lcoh$ while leaving $\lvis$ near zero.

Applying the noise once per layer accumulates the aligned contraction into the layered coherence defect $\Dcohlay(\gamma)=\gamma L\,\lvis+O((\gamma L)^2)$, which reduces to $\gamma L\,\lcoh$ in the restricted-isotropic case. For the four restricted-isotropic channels used in the main sweep, the worst-case and aligned rates agree to four significant figures, with $\lcoh=1.000$ for amplitude damping, $3.996$ for dephasing, $6.648$ for depolarising and $7.973$ for the $U(1)$-breaking $X$-error control. Here, restricted-isotropic refers to the action on the intra-sector off-diagonal block rather than to ordinary channel isotropy.


\begin{theorem}[Perturbative coherence-loss law]\label{thm:perturb}
For an active parameter $\theta_i$ and a phase-covariant Markovian noise channel of rate $\gamma$ with Frobenius-dissipative restricted generator acting once per layer \cite{holevo1996,lindblad1976,gks1976}, and under the first-order alignment condition that the noisy derivative acquires no leading-order component along readout-visible directions orthogonal to $G_i$, the leading-order degradation is a sum of layer-resolved aligned contractions,
\begin{equation}
    \Mtwo(\theta_i;\gamma)
    =
    \Mtwo(\theta_i;0)
    \Bigl[1-2\,\gamma\textstyle\sum_{\ell=1}^{L}\lvis^{(\ell)}(\theta_i)+O\bigl((\gamma L\,\lcoh)^2\bigr)\Bigr].
    \label{eq:perturb}
\end{equation}
In the restricted-isotropic family, where every off-diagonal direction decays at the common rate, the layer sum collapses to a scalar accumulation,
\begin{equation}
    \Mtwo(\theta_i;\gamma)
    =
    \Mtwo(\theta_i;0)
    \bigl[1-2\,\gamma L\,\lvis(\theta_i)+O\bigl((\gamma L\,\lcoh)^2\bigr)\bigr],
    \label{eq:perturbscalar}
\end{equation}
with the bounded alignment factor $\omega_i=\lvis(\theta_i)/\lcoh\in[0,1]$ equal to one for isotropic generators and strictly less than one when the gradient mode is misaligned with the fastest-decaying direction.
\end{theorem}


Because three closely related alignment quantities appear in the paper, Table~\ref{tab:alignment} collects their definitions. The theoretical ratio $\omega_i=\lvis/\lcoh$ is the object the theory identifies as fundamental, its perturbative empirical estimate $\weff$ is read off the small-noise experiments, and the finite-noise multiplier $\aeff$ is an a posteriori diagnostic measured in the presence of finite noise.

\begin{table}[t]
\caption{The three alignment quantities used in the paper. The first is a theoretical ratio, the second its perturbative empirical estimate, and the third a finite-noise empirical multiplier.}\label{tab:alignment}
\centering
\begin{tabular}{@{}lll@{}}
\toprule
Symbol & Definition & Role\\
\midrule
$\omega_i=\lvis/\lcoh$ & ratio of aligned to worst-case rate & theoretical alignment factor, in $[0,1]$\\
$\weff$ & $\widetilde{\Delta}/(2\gamma L\,\lcoh)$ in the perturbative window & empirical estimate of $\omega_i$ at small noise\\
$\aeff$ & finite-noise weight fitted to paired degradation & a posteriori diagnostic of residual alignment\\
\bottomrule
\end{tabular}
\end{table}


\begin{proof}[Proof of Theorem~\ref{thm:perturb}]
Let $\Phi^\gamma=I+\gamma\Lop+O(\gamma^2)$ be the small-noise expansion of the layer channel and $\Pioff$ the projector onto the off-diagonal block of sector $r$. The worst-case rate is Equation~\eqref{eq:lcohdef} and the aligned rate is the Rayleigh quotient of Equation~\eqref{eq:lvisdef} with respect to the gradient mode $G_i$. Inserting the expansion into the commutator form of the derivative and projecting onto $G_i(\theta)=\Pioff[Z_{q_i},O_B(\ell_i;\theta)]$ gives, at each layer at which the noise acts, a first-order contraction of the gradient mode propagated to that layer at the pointwise rate $\lvis^{(\ell)}(\theta)$. The single-layer contraction along the mode is the Rayleigh quotient by definition, and for a Frobenius-dissipative restricted generator, each such quotient is a non-negative spectral decay rate. Because the noise acts once per layer and the first-order terms add, the leading-order degradation is the layer sum $2\gamma\sum_{\ell=1}^{L}\lvis^{(\ell)}(\theta_i)$ of Equation~\eqref{eq:perturb} after the $\theta$ ensemble average.

Writing the first-order noisy derivative as $\partial_{\theta_i}f^\gamma_\theta=\partial_{\theta_i}f^0_\theta(1-\gamma\sum_\ell\lvis^{(\ell)})+\gamma\,\delta_i$, where $\delta_i$ collects the first-order response along readout-visible directions orthogonal to $G_i$, the square produces a diagonal term and cross terms linear in $\delta_i$. The first-order alignment condition is exactly the statement that the ensemble average of these cross terms vanishes at leading order, which holds identically in the restricted-isotropic family where $\delta_i=0$ because every off-diagonal direction decays at the common rate. The collapse to the scalar accumulation $\gamma L\,\lvis$ of Equation~\eqref{eq:perturbscalar} requires the aligned rate to be layer-independent, which holds in that family because the interleaved unitaries rotate the gradient mode without changing its decay rate. Outside it, the interleaved unitaries can carry the mode through directions of differing aligned rate, the layer-resolved rates no longer coincide, and the scalar $\lcoh$ can no longer stand in for the layer-averaged aligned contraction. The alignment ratio $\omega_i=\lvis/\lcoh$ inherits the bounds $[0,1]$ from Equation~\eqref{eq:visbound}.
\end{proof}


\begin{thebibliography}{99}
\bibitem[breuerpetruccione2002]{breuerpetruccione2002} H.-P.~Breuer and F.~Petruccione, \emph{The Theory of Open Quantum Systems}. Oxford University Press, Oxford, 2002. \url{https://doi.org/10.1093/acprof:oso/9780199213900.001.0001}
\bibitem[gks1976]{gks1976} V.~Gorini, A.~Kossakowski, and E.~C.~G.~Sudarshan, ``Completely positive dynamical semigroups of $N$-level systems,'' \emph{Journal of Mathematical Physics}, vol.~17, pp.~821--825, 1976. \url{https://doi.org/10.1063/1.522979}
\bibitem[grant2019]{grant2019} E.~Grant, L.~Wossnig, M.~Ostaszewski, and M.~Benedetti, ``An initialization strategy for addressing barren plateaus in parametrized quantum circuits,'' \emph{Quantum}, vol.~3, art.~214, 2019. \url{https://doi.org/10.22331/q-2019-12-09-214}
\bibitem[holevo1996]{holevo1996} A.~S.~Holevo, ``Covariant quantum Markovian evolutions,'' \emph{Journal of Mathematical Physics}, vol.~37, pp.~1812--1832, 1996. \url{https://doi.org/10.1063/1.531481}
\bibitem[larocca2022groupinv]{larocca2022groupinv} M.~Larocca \emph{et al.}, ``Group-invariant quantum machine learning,'' \emph{PRX Quantum}, vol.~3, art.~030341, 2022. \url{https://doi.org/10.1103/PRXQuantum.3.030341}
\bibitem[lindblad1976]{lindblad1976} G.~Lindblad, ``On the generators of quantum dynamical semigroups,'' \emph{Communications in Mathematical Physics}, vol.~48, pp.~119--130, 1976. \url{https://doi.org/10.1007/BF01608499}
\bibitem[mitarai2018]{mitarai2018} K.~Mitarai, M.~Negoro, M.~Kitagawa, and K.~Fujii, ``Quantum circuit learning,'' \emph{Physical Review A}, vol.~98, art.~032309, 2018. \url{https://doi.org/10.1103/PhysRevA.98.032309}
\bibitem[nguyen2024]{nguyen2024} Q.~T.~Nguyen \emph{et al.}, ``Theory for equivariant quantum neural networks,'' \emph{PRX Quantum}, vol.~5, art.~020328, 2024. \url{https://doi.org/10.1103/PRXQuantum.5.020328}
\bibitem[schatzki2024]{schatzki2024} L.~Sch\"atzki, M.~Larocca, Q.~T.~Nguyen, F.~Sauvage, and M.~Cerezo, ``Theoretical guarantees for permutation-equivariant quantum neural networks,'' \emph{npj Quantum Information}, vol.~10, art.~12, 2024. \url{https://doi.org/10.1038/s41534-024-00804-1}
\bibitem[schuld2019]{schuld2019} M.~Schuld, V.~Bergholm, C.~Gogolin, J.~Izaac, and N.~Killoran, ``Evaluating analytic gradients on quantum hardware,'' \emph{Physical Review A}, vol.~99, art.~032331, 2019. \url{https://doi.org/10.1103/PhysRevA.99.032331}
\bibitem[wierichs2022]{wierichs2022} D.~Wierichs, J.~Izaac, C.~Wang, and C.~Y.-Y.~Lin, ``General parameter-shift rules for quantum gradients,'' \emph{Quantum}, vol.~6, art.~677, 2022. \url{https://doi.org/10.22331/q-2022-03-30-677}
\end{thebibliography}
\end{document}
